\section{Exam map and control concepts}

\subsection{Ctrl+F keywords}

\subsubsection{Course keywords, eksamen, regulering, feedback, controller}
\[
\{r,e,u,y,y_m,G,C,H,L,S,T,\omega_c,\gamma_M,A_M,e_{\mathrm{ss}}\}
\]
Use this file as the front-door map. Danish search aliases used below include regulering, tilbagekobling, forstyrrelse, maaling, folgefejl, stabilitet, and bodeplot.

\subsection{Signal and loop notation}

\subsubsection{Standard feedback signals}
\[
e(t)=r(t)-y_m(t),\qquad y_m(t)=H(s)y(t),\qquad u(t)=C(s)e(t)
\]
\(r\) is reference, \(e\) is error, \(u\) is plant input, \(y\) is output, and \(y_m\) is measured output.

\subsubsection{Plant, controller, measurement, loop transfer}
\[
L(s)=C(s)G(s)H(s)
\]
Most stability-margin and sensitivity questions start by identifying the correct loop transfer function \(L\).

\subsubsection{Closed-loop transfer and characteristic equation}
\[
T(s)=\frac{L(s)}{1+L(s)},\qquad 1+L(s)=0
\]
Closed-loop stability is decided by the roots of the characteristic equation, not by open-loop poles alone.

\subsection{Exam problem recognition}

\subsubsection{Problem type keyword table}
\begin{center}
\begin{tabular}{>{\raggedright\arraybackslash}p{0.38\linewidth}|>{\raggedright\arraybackslash}p{0.50\linewidth}}
\textbf{Exam wording} & \textbf{Go to}\\ \hline
block diagram, feedback, summing junction & Block diagrams and feedback\\
RLC, masses, equations of motion, model & Modelling and linearization\\
differential equation, Laplace, poles, zeros & Laplace transform and transfer functions\\
step response, overshoot, steady-state value & Frequency and time-domain analysis\\
Bode, gain margin, phase margin, \(K_P\) & Bode plots and stability margins\\
Nyquist, encirclement, unstable open loop & Nyquist plot and unstable systems\\
PI-Lead, Lead, Lag, \(\alpha,\beta,N_i\) & PI-Lead design and limited systems\\
disturbance, sensitivity, prefilter, feed-forward & Disturbances and feed-forward
\end{tabular}
\end{center}
Use the wording in the problem statement first, then confirm with the formulas in the relevant section.

\subsection{Controller templates}

\subsubsection{P, PI, Lead, Lag controller forms}
\[
\begin{aligned}
C_P(s)&=K_P,\\
C_{PI}(s)&=K_P\frac{\tau_i s+1}{\tau_i s},\\
C_{\mathrm{lead}}(s)&=\frac{\tau_d s+1}{\alpha\tau_d s+1},\quad 0<\alpha<1,\\
C_{\mathrm{lag}}(s)&=\frac{\tau_i s+1}{\frac{\tau_i}{\beta}s+1},\quad \beta>1.
\end{aligned}
\]
These are the recurring controller blocks in the lectures and F21 exam.

\subsubsection{Performance and robustness symbols}
\[
\begin{array}{c|l}
t_r & \text{rise time}\\
t_s & \text{settling time}\\
M_p & \text{percent overshoot}\\
e_{\mathrm{ss}} & \text{steady-state error, stationaer fejl}\\
\omega_c & \text{gain crossover frequency}\\
\gamma_M & \text{phase margin}\\
A_M & \text{gain margin}
\end{array}
\]
Fast exam checks often use signs, units, limits, and whether these specifications move in the expected direction.

\subsection{Multiple-choice and fast checks}

\subsubsection{Elimination rules for answer options}
\[
\text{wrong sign} \;\vee\; \text{wrong units} \;\vee\; \text{unstable hidden pole} \;\vee\; \text{improper controller}
\quad\Longrightarrow\quad \text{reject option}
\]
For multi-clause answers, one false clause makes the option false.

\subsubsection{Limit checks}
\[
\lim_{K_P\to 0}T(s)=0,\qquad \lim_{K_P\to\infty}\frac{K_PG(s)}{1+K_PG(s)}\approx 1
\]
Use limiting cases to reject transfer functions with impossible signs or missing feedback denominators.
